\section{Roles y Usuarios}
\subsection{Módulo de Role-Users (/roleusers)}

\subsection{Asignar Rol a Usuario (POST /roleusers)}

\casodeprueba{CP1.1}
{Asignar rol a usuario ok}
{Token JWT con permiso ROLEUSERS\_CREATE, usuario existente y rol existente.}
{JSON con datos válidos { 'id\_user': 'f47ac10b-58cc-4372-a567-0e02b2c3d479', 'role\_name': 'ADMIN' }.}
{\begin{enumerate} [leftmargin=*, nosep]
    \item Código de respuesta 201
    \item Respuesta con datos correctos de la asignación creada.
    \item Buscar asignación específica (GET /roleusers/:id\_user/:id\_role) para confirmar creación.
\end{enumerate}}

\casodeprueba{CP1.2}
{Asignar rol a usuario duplicado}
{Token JWT con permisos suficientes y asignación ya existente para el usuario y rol.}
{JSON con id\_user y role\_name ya asignados previamente.}
{\begin{enumerate} [leftmargin=*, nosep]
    \item Código de respuesta 409
    \item Mensaje de error indicando que la asignación ya existe.
\end{enumerate}}

\casodeprueba{CP1.3.1}
{Asignar rol - id\_user vacío o ausente}
{Token JWT con permisos suficientes.}
{JSON con id\_user vacío ('') o ausente.}
{\begin{enumerate} [leftmargin=*, nosep]
    \item Código de respuesta 400
    \item Mensaje de error indicando que el id de usuario es obligatorio.
\end{enumerate}}

\casodeprueba{CP1.3.2}
{Asignar rol - id\_user no es UUID v4 válido}
{Token JWT con permisos suficientes.}
{JSON con id\_user de formato inválido (ej. '123-abc').}
{\begin{enumerate} [leftmargin=*, nosep]
    \item Código de respuesta 400
    \item Mensaje de error indicando que id\_user debe ser un UUID v4 válido.
\end{enumerate}}

\casodeprueba{CP1.3.3}
{Asignar rol - role\_name vacío o ausente}
{Token JWT con permisos suficientes.}
{JSON con role\_name vacío ('') o ausente.}
{\begin{enumerate} [leftmargin=*, nosep]
    \item Código de respuesta 400
    \item Mensaje de error indicando que el nombre del rol es obligatorio.
\end{enumerate}}

\casodeprueba{CP1.3.4}
{Asignar rol - role\_name no es string}
{Token JWT con permisos suficientes.}
{JSON con role\_name numérico (123).}
{\begin{enumerate} [leftmargin=*, nosep]
    \item Código de respuesta 400
    \item Mensaje de error indicando que el nombre del rol debe ser una cadena de texto.
\end{enumerate}}

\casodeprueba{CP2.1}
{Listar asignaciones rol-usuario ok}
{Token JWT con permiso ROLEUSERS\_READ y asignaciones registradas.}
{Ninguna (GET /roleusers).}
{\begin{enumerate} [leftmargin=*, nosep]
    \item Código de respuesta 200
    \item Cantidad de registros obtenidos en el array.
\end{enumerate}}

\casodeprueba{CP2.2}
{Listar asignaciones sin resultados}
{Token JWT con permiso ROLEUSERS\_READ y base de datos sin asignaciones.}
{Ninguna (GET /roleusers).}
{\begin{enumerate} [leftmargin=*, nosep]
    \item Código de respuesta 200
    \item Cantidad de registros (0).
\end{enumerate}}

\casodeprueba{CP3.1}
{Obtener roles por usuario ok}
{Token JWT con permiso ROLEUSERS\_READ y usuario con roles asignados.}
{UUID válido del usuario como parámetro (GET /roleusers/user/:id\_user).}
{\begin{enumerate} [leftmargin=*, nosep]
    \item Código de respuesta 200
    \item Asignaciones del usuario obtenidas correctamente.
\end{enumerate}}

\casodeprueba{CP3.2}
{Obtener roles por usuario - Usuario not found}
{Token JWT con permiso ROLEUSERS\_READ.}
{UUID de usuario inexistente.}
{\begin{enumerate} [leftmargin=*, nosep]
    \item Código de respuesta 404
    \item Mensaje de error de usuario no encontrado.
\end{enumerate}}

\casodeprueba{CP3.3}
{Obtener roles por usuario - UUID inválido}
{Token JWT con permiso ROLEUSERS\_READ.}
{UUID con formato inválido.}
{\begin{enumerate} [leftmargin=*, nosep]
    \item Código de respuesta 400
    \item Mensaje de error.
\end{enumerate}}

\casodeprueba{CP4.1}
{Obtener usuarios por nombre de rol ok}
{Token JWT con permiso ROLEUSERS\_READ, rol existente con usuarios asignados.}
{Nombre del rol como parámetro (GET /roleusers/role/ADMIN).}
{\begin{enumerate} [leftmargin=*, nosep]
    \item Código de respuesta 200
    \item Lista de usuarios con el rol obtenida correctamente.
\end{enumerate}}

\casodeprueba{CP4.2}
{Obtener usuarios por rol - Rol not found / sin resultados}
{Token JWT con permiso ROLEUSERS\_READ.}
{Nombre de rol inexistente.}
{\begin{enumerate} [leftmargin=*, nosep]
    \item Código de respuesta 404 (o 200 con 0 registros)
    \item Mensaje de error o array vacío.
\end{enumerate}}

\casodeprueba{CP5.1}
{Obtener asignación específica ok}
{Token JWT con permiso ROLEUSERS\_READ y asignación existente.}
{id\_user (UUID) y id\_role (UUID) como parámetros.}
{\begin{enumerate} [leftmargin=*, nosep]
    \item Código de respuesta 200
    \item Datos de la respuesta coinciden con la asignación buscada.
\end{enumerate}}

\casodeprueba{CP5.2}
{Obtener asignación específica not found}
{Token JWT con permiso ROLEUSERS\_READ.}
{id\_user y id\_role válidos pero sin asignación mutua.}
{\begin{enumerate} [leftmargin=*, nosep]
    \item Código de respuesta 404
    \item Mensaje de error de asignación no encontrada.
\end{enumerate}}

\casodeprueba{CP5.3}
{Obtener asignación específica - UUID inválido}
{Token JWT con permiso ROLEUSERS\_READ.}
{Parámetros con formato no UUID.}
{\begin{enumerate} [leftmargin=*, nosep]
    \item Código de respuesta 400
    \item Mensaje de error.
\end{enumerate}}

\casodeprueba{CP6.1}
{Activar asignación ok}
{Token JWT con permiso ROLEUSERS\_UPDATE y asignación desactivada.}
{id\_user y id\_role válidos como parámetros.}
{\begin{enumerate} [leftmargin=*, nosep]
    \item Código de respuesta 200
    \item Asignación activada exitosamente.
\end{enumerate}}

\casodeprueba{CP6.2}
{Activar asignación not found}
{Token JWT con permiso ROLEUSERS\_UPDATE.}
{IDs sin asignación existente.}
{\begin{enumerate} [leftmargin=*, nosep]
    \item Código de respuesta 404
    \item Mensaje de error.
\end{enumerate}}

\casodeprueba{CP7.1}
{Desactivar asignación ok}
{Token JWT con permiso ROLEUSERS\_DELETE y asignación activa.}
{id\_user y id\_role válidos como parámetros.}
{\begin{enumerate} [leftmargin=*, nosep]
    \item Código de respuesta 200
    \item Asignación desactivada exitosamente.
    \item Buscar asignación para confirmar cambio de estado.
\end{enumerate}}

\casodeprueba{CP7.2}
{Desactivar asignación not found}
{Token JWT con permiso ROLEUSERS\_DELETE.}
{IDs sin asignación existente.}
{\begin{enumerate} [leftmargin=*, nosep]
    \item Código de respuesta 404
    \item Mensaje de error.
\end{enumerate}}

\casodeprueba{CP8.1}
{Eliminar asignación ok}
{Token JWT con permiso ROLEUSERS\_DELETE y asignación existente.}
{id\_user y id\_role válidos como parámetros.}
{\begin{enumerate} [leftmargin=*, nosep]
    \item Código de respuesta 200
    \item Asignación eliminada exitosamente.
    \item Buscar asignación (GET /roleusers/:id\_user/:id\_role) para confirmar not found (404).
\end{enumerate}}

\casodeprueba{CP8.2}
{Eliminar asignación not found}
{Token JWT con permiso ROLEUSERS\_DELETE.}
{IDs sin asignación existente.}
{\begin{enumerate} [leftmargin=*, nosep]
    \item Código de respuesta 404
    \item Mensaje de error.
\end{enumerate}}

